\documentclass[11pt,a4paper]{article}

% --- Encodage / langue ---
\usepackage[utf8]{inputenc}
\usepackage[T1]{fontenc}
\usepackage[french]{babel}
\usepackage{lmodern}

% --- Mise en page compacte ---
\usepackage[a4paper,margin=2cm]{geometry}
\usepackage{microtype}
\usepackage{enumitem}
\setlist{leftmargin=1.4em}
\usepackage{titlesec}
\titlespacing*{\section}{0pt}{8pt}{2pt}
\titleformat{\section}{\normalfont\large\bfseries}{\thesection.}{0.5em}{}
\setlength{\parindent}{0pt}
\setlength{\parskip}{2pt}

\usepackage[backend=biber,style=numeric-comp,sorting=none,maxbibnames=3]{biblatex}
\addbibresource{references.bib}
\renewcommand*{\bibfont}{\footnotesize}

\usepackage{hyperref}
\hypersetup{hidelinks}

\title{\vspace{-1.2cm}\textbf{Livrable 2 --- Base d'évaluation\\[2pt]
\large Analyse automatique de partitions musicales (OMR)}}

\author{\normalsize Thibault Bourdoiseau \& Léo Li --- SAE TIFO 1 / TIFO 2 / MLRF}
\date{\vspace{-8pt}\small \today}

\begin{document}
\maketitle
\vspace{-10pt}

Le sujet laisse la base libre. Nous combinons deux bases publiques déjà annotées --- une \textbf{imprimée} et
une \textbf{manuscrite}, pour comparer la difficulté des deux cas --- complétées par une petite annotation
manuelle pour les hauteurs et durées.

\section{DeepScoresV2 --- imprimé~\cite{tuggener2020}}
\textbf{Description.} Grande base de partitions imprimées (typographiées) annotées pour la détection
d'objets~: pour chaque image, les \textbf{boîtes englobantes étiquetées} de $135$ types de symboles musicaux
(têtes de notes, clefs, silences, altérations\dots). L'annotation est livrée dans un fichier JSON séparé
(format type COCO). Elle alimente le critère de \emph{F1 de détection des têtes de notes}. On en prélève un
sous-ensemble de quelques dizaines d'images, découpé par niveau de difficulté (police propre, bruit,
déformation).

\textbf{Lien.} \url{https://zenodo.org/records/4012193}

\textbf{Autorisation d'accès.} Licence \textbf{Creative Commons CC BY 4.0}~: réutilisation libre, y compris la
redistribution, à condition de citer Tuggener et al.~(2020). Téléchargement direct sur Zenodo, sans
inscription.

\section{MUSCIMA++ --- manuscrit~\cite{hajic2017,fornes2012}}
\textbf{Description.} $140$ pages de partitions \textbf{manuscrites}, annotées au niveau du symbole~:
$91\,000$ symboles avec leurs \textbf{boîtes englobantes et leurs classes} (et les relations entre symboles).
C'est l'équivalent manuscrit de DeepScoresV2~; il alimente le même critère de \emph{F1 de détection}, sur des
images bien plus difficiles (portées irrégulières, têtes déformées), ce qui permet de mesurer l'écart de
performance imprimé~/~manuscrit. Les images sont un sous-ensemble de CVC-MUSCIMA~\cite{fornes2012}.

\textbf{Lien.} \url{https://ufal.mff.cuni.cz/muscima}

\textbf{Autorisation d'accès.} Licence \textbf{Creative Commons CC BY-NC-SA 4.0}~: usage recherche / non
commercial, partage à l'identique, avec citation de Haji\v{c} jr.\ \& Pecina~(2017) et de Fornés et
al.~(2012).

\section{Annotation manuelle des hauteurs et durées (PA / RA)}
Les critères de \emph{précision des hauteurs} et de \emph{durées} exigent une vérité-terrain donnant la
\textbf{mélodie décodée} (suite de notes \texttt{hauteur~+~durée}). Cette information n'est fournie ni par
DeepScoresV2 ni par MUSCIMA++, qui ne donnent que la position et la classe des symboles~; la déduire des
annotations reviendrait à réimplémenter le décodage que le système doit justement réaliser (évaluation
circulaire). Nous constituons donc à la main une \textbf{petite base} de quelques images (prises dans les deux
bases ci-dessus, imprimées et manuscrites), en saisissant pour chacune la séquence de notes attendue au format
JSON des outils (\texttt{\{"notes": [\{"pitch": "C4", "duration": "q"\}, \dots]\}}). C'est peu coûteux (lecture
et saisie, sans tracé de boîtes) et conforme à la recommandation du sujet d'une \emph{petite} base
d'évaluation.

\printbibliography[heading=bibnumbered,title={Références}]

\end{document}
